% =====================================================================
%  5.13 ESPECIFICACION Y ANALISIS DE REQUISITOS
%  Rubrica (Doc2): requisitos Y analisis, ambos "descritos con detalle".
%  Nota: el 30% de la nota de Desarrollo es "Cumprimento dos requisitos",
%  asi que los requisitos que escribas aqui son la vara de medir con la
%  que se evaluara tu software. Escribelos con cuidado.
% =====================================================================

\chapter{Especificación y análisis de requisitos}
\label{cap:requisitos}

Los requisitos que siguen describen el sistema \emph{tal como está implementado}.
No se incluyen capacidades deliberadamente fuera de alcance
(\cref{sec:limitaciones}): integración con la historia clínica electrónica,
paso dedicado de ficha demográfica ni sustitución del juicio clínico. Cada
requisito funcional es verificable sobre la interfaz o sobre el resultado de
la evaluación. Los criterios de aceptación observables se recogen en
el~\cref{anexo:aceptacion} (\cref{tab:aceptacion-rf}); no presuponen que la
prueba se haya ejecutado ni sustituyen la evidencia del~\cref{cap:pruebas}.

\section{Actores del sistema}
\label{sec:actores}

El actor primario es el \textbf{profesional sanitario} que revisa la
medicación de una persona mayor ---geriatra, médico clínico o farmacéutico
asistencial---. Interactúa con la aplicación en consulta: selecciona fármacos
y diagnósticos, anota analítica cuando procede y examina los avisos STOPP y
START. El sistema es de \emph{apoyo a la decisión}: muestra criterios
aplicables, pero no prescribe, no cierra el caso clínico y no sustituye la
interpretación del profesional~\cite{omahony2023stopp}.

No hay actores secundarios de sistema (servidor, historia clínica, laboratorio
externo). La persistencia y la evaluación ocurren en el navegador del propio
profesional, según la visión del~\cref{cap:solucion}.

\section{Requisitos funcionales}
\label{sec:rf}

La~\cref{tab:rf} recoge los requisitos funcionales con identificador estable,
prioridad y descripción sin ambigüedad. La prioridad \emph{Alta} corresponde a
capacidades sin las cuales el flujo de revisión no se completa; \emph{Media},
a ayudas de captura, accesibilidad o reutilización del caso.

{\small
\begin{longtable}{@{}llp{9.2cm}@{}}
  \caption{Requisitos funcionales del sistema implementado.}
  \label{tab:rf} \\
  \toprule
  \textbf{ID} & \textbf{Prioridad} & \textbf{Descripción} \\
  \midrule
  \endfirsthead
  \caption[]{Requisitos funcionales del sistema implementado (cont.).} \\
  \toprule
  \textbf{ID} & \textbf{Prioridad} & \textbf{Descripción} \\
  \midrule
  \endhead
  \bottomrule
  \endfoot
  RF01 & Alta & El profesional podrá seleccionar y deseleccionar medicamentos
    del catálogo, organizados por categoría terapéutica (pestañas
    cardiovascular, anticoagulantes, SNC, renal, gastrointestinal,
    respiratorio, endocrino, urológico, osteoarticular, antiinfecciosos y
    «Otros»). Cada fármaco marcado formará parte del caso y podrá
    desmarcarse en cualquier momento. \\
  RF02 & Alta & El profesional podrá seleccionar y deseleccionar diagnósticos
    del catálogo, organizados por sistema. En las familias con variantes
    mutuamente excluyentes ---en la implementación actual, la hipertensión
    arterial: HTA sin especificar, no complicada, moderada o grave---
    marcar una variante retirará las hermanas. Los diagnósticos que
    dependen de una medicación habilitante permanecerán no seleccionables
    hasta que dicha medicación esté en el caso. \\
  RF03 & Alta & En la pestaña «Otros» del paso de diagnósticos, el profesional
    podrá introducir, modificar o dejar vacíos los parámetros de analítica y
    constantes que algún criterio lee: TFGe, TSH, frecuencia cardíaca, QTc,
    potasio, sodio, calcio corregido, PAS y PAD. El panel será fijo (no
    condicionado a la selección). \\
  RF04 & Alta & Tras cada cambio de medicación, diagnóstico o analítica, el
    sistema evaluará automáticamente el caso frente al catálogo STOPP/START
    y actualizará los avisos aplicables en el propio paso, sin petición
    explícita ni ida y vuelta a un servidor. \\
  RF05 & Baja & El motor ofrecerá un cálculo auxiliar de los medicamentos
    del catálogo que, de añadirse al caso actual, dispararían una regla con
    campo \texttt{excludes}. Esta capacidad no se presenta en la interfaz
    actual y queda disponible para una integración posterior. \\
  RF06 & Alta & En la pestaña activa, el sistema mostrará un bloque
    «Relevantes de otros sistemas» con fármacos o diagnósticos de otras
    pestañas que participan en criterios del sistema en revisión, de modo
    que el profesional no deba abandonar la pestaña para marcarlos. \\
  RF07 & Media & El profesional podrá marcar una pestaña como revisada cuando
    no tenga selección, para dejar constancia de que la ha examinado y no
    aplica nada. Si añade una selección, el marcador se retirará. Una
    pestaña con selección no podrá marcarse como revisada. \\
  RF08 & Alta & El sistema mostrará los criterios aplicables agrupados por
    tipo (STOPP y START) y, dentro de cada tipo, por sistema orgánico,
    con código de sección y texto del aviso. \\
  RF09 & Alta & El profesional podrá descargar un informe PDF del caso con
    cabecera institucional, fecha, diagnósticos, medicaciones y tablas de
    criterios STOPP y START aplicables. \\
  RF10 & Media & El profesional podrá copiar al portapapeles un texto plano
    con los criterios aplicables, agrupados por sistema. Si no hay
    avisos, el texto será «Sin criterios aplicables.». \\
  RF11 & Alta & El profesional podrá exportar el caso completo a un fichero
    JSON versionado e importar un fichero del mismo formato. Un JSON que
    no supere la validación se rechazará sin alterar el caso en curso. \\
  RF12 & Media & El profesional podrá reiniciar el caso. El sistema pedirá
    confirmación explícita; solo si se confirma se vaciarán medicaciones,
    diagnósticos, analítica y marcadores de revisión. \\
  RF13 & Media & El profesional podrá elegir la escala de fuente (pequeño,
    mediano o grande) y la orientación de las pestañas (raíl vertical u
    horizontal). Ambas preferencias se conservarán entre recargas. \\
  RF14 & Media & El profesional podrá abrir una guía rápida que explique STOPP
    y START, el flujo en dos pasos y el enlace al PDF oficial de la
    guía~\cite{sacyl2023stopp}. \\
  RF15 & Media & El profesional podrá añadir, en un grupo del catálogo, un
    ítem «Otro» que no figura como fármaco o diagnóstico nominado, de modo
    que la clase o el grupo participe en la evaluación aunque el principio
    activo concreto no esté listado. \\
\end{longtable}
}

\subsection{Criterios de aceptación de los requisitos funcionales}

Para evitar repetir cada descripción, la~\cref{tab:aceptacion-rf} del
\cref{anexo:aceptacion} agrupa requisitos que comparten procedimiento. La
comprobación parte de un caso vacío salvo que el propio escenario requiera
datos previos.

\section{Requisitos no funcionales}
\label{sec:rnf}

La~\cref{tab:rnf} fija cinco condiciones de calidad, cada una con un
criterio observable. Se toman como referencia las características de
ISO/IEC~25010~\cite{iso25010}: usabilidad del flujo, ejecución en el
cliente, adaptación visual, mantenibilidad de las reglas y tratamiento
local del caso. No se incluye un tiempo máximo de respuesta porque no se
realizó ninguna medición.

\begin{table}[H]
  \centering
  \scriptsize
  \begin{tabularx}{\textwidth}{@{}l>{\raggedright\arraybackslash}p{2.3cm}XX@{}}
    \toprule
    \textbf{ID} & \textbf{Característica} & \textbf{Requisito} &
      \textbf{Criterio de aceptación} \\
    \midrule
    RNF01 & Usabilidad & El flujo tendrá dos pasos
      (medicamentos y diagnósticos) y mostrará los avisos en la misma
      vista que la captura. &
      Se recorren las dos rutas y en ambas se puede seleccionar una
      entrada y observar sus avisos sin abrir una tercera pantalla. \\
    RNF02 & Eficiencia de desempeño & La evaluación de criterios se
      ejecutará en el cliente. Después de cargar los activos, un cambio
      del caso no requerirá una petición de red del motor. &
      Con el catálogo ya cargado, una prueba cambia una entrada y
      observa la actualización; la inspección de red no muestra una
      petición nueva asociada a evaluar el caso. No se evalúa latencia. \\
    RNF03 & Adaptabilidad visual & Existirán tres escalas
      de fuente ($1$, $1{,}15$ y $1{,}3$) aplicables a toda la interfaz,
      con persistencia de la elección. &
      Cada opción cambia la variable de escala global y el valor elegido
      permanece tras recargar. Este requisito no acredita por sí solo
      conformidad WCAG ni usabilidad clínica. \\
    RNF04 & Mantenibilidad & Las reglas clínicas residirán en un catálogo
      JSON independiente del código de la interfaz, de forma que un
      criterio pueda corregirse o añadirse sin reescribir el motor. &
      El motor carga una regla válida desde el catálogo o una prueba y la
      evalúa sin una rama específica para su identificador; los operadores
      de dominio siguen siendo código cuando la expresión los necesita. \\
    RNF05 & Seguridad / privacidad & El caso no se enviará a un servidor
      de aplicación; la persistencia será local y el reinicio permitirá
      retirar los valores clínicos del estado persistido
      (\cref{cap:datos}). &
      La inspección de red durante captura y evaluación no contiene el
      caso; tras confirmar el reinicio el estado restaurable es un caso
      vacío, aunque puedan persistir claves técnicas con valores vacíos.
      Un JSON de importación inválido se rechaza sin sustituir el caso. \\
    \bottomrule
  \end{tabularx}
  \caption{Requisitos no funcionales, referidos al modelo de calidad
  ISO/IEC~25010.}
  \label{tab:rnf}
\end{table}

\section{Casos de uso}
\label{sec:casos-uso}

El actor principal es el profesional sanitario. Los casos de uso describen
sus objetivos sobre el sistema, tratado como caja negra: no importan las
clases internas, sino lo que el profesional puede hacer. El caso de uso
principal es revisar la medicación de un paciente con apoyo STOPP/START.
De él se desprenden la captura (RF01--RF03, RF15), la consulta de avisos
(RF08) y las salidas (RF09--RF11). Dos subcasos ocurren siempre que cambia
el caso, sin que el profesional los invoque aparte: \textbf{evaluar
criterios} (RF04) y \textbf{persistir el caso} en el navegador. La
exportación JSON sí es una acción explícita. La~\cref{fig:casos-uso}
resume estas relaciones.

\begin{figure}[H]
  \centering
  \resizebox{\textwidth}{!}{%
  \begin{tikzpicture}[
    font=\small\sffamily,
    uc/.style={
      ellipse, draw, thick, align=center,
      minimum width=3.55cm, minimum height=0.92cm,
      inner sep=1pt, fill=gray!6
    },
    ucinc/.style={
      ellipse, draw, thick, align=center,
      minimum width=3.15cm, minimum height=0.92cm,
      inner sep=1pt, fill=gray!16
    },
    assoc/.style={thick},
    inc/.style={-{Latex}, thick, dashed}
  ]
    % System boundary (origin)
    \node[
      draw, thick, inner sep=0pt,
      minimum width=11.6cm, minimum height=10.6cm
    ] (sys) {};
    \node[font=\small\bfseries\sffamily, fill=white, inner sep=3pt]
      at (sys.north) {Sistema STOPP/START};

    % Actor, vertically centered with the system
    \node[inner sep=0pt, minimum width=0.7cm, minimum height=1.15cm,
      left=1.35cm of sys] (act) {};
    \begin{scope}[shift={(act.center)}]
      \fill (0,0.40) circle (0.11);
      \draw[thick, line cap=round]
        (0,0.29) -- (0,-0.10)
        (-0.20,0.16) -- (0.20,0.16)
        (0,-0.10) -- (-0.16,-0.42)
        (0,-0.10) -- (0.16,-0.42);
    \end{scope}
    \node[below=0.08cm of act, align=center, font=\small\bfseries\sffamily]
      (actlbl) {Profesional\\[-1pt]{\normalfont\scriptsize sanitario}};

    % Primary use cases (left column, actor-facing)
    \node[uc] at ([shift={(-2.95, 4.45)}]sys.center) (ucmed)
      {Seleccionar\\medicación};
    \node[uc] at ([shift={(-2.95, 3.20)}]sys.center) (ucdx)
      {Seleccionar\\diagnósticos};
    \node[uc] at ([shift={(-2.95, 1.95)}]sys.center) (uclab)
      {Introducir\\analítica};
    \node[uc] at ([shift={(-2.95, 0.70)}]sys.center) (ucver)
      {Consultar avisos\\STOPP / START};
    \node[uc] at ([shift={(-2.95,-0.55)}]sys.center) (ucexp)
      {Exportar informe};
    \node[uc] at ([shift={(-2.95,-1.80)}]sys.center) (ucges)
      {Importar o\\reiniciar el caso};
    \node[uc] at ([shift={(-2.95,-3.05)}]sys.center) (ucvis)
      {Ajustar\\visualización};
    \node[uc] at ([shift={(-2.95,-4.30)}]sys.center) (ucguia)
      {Consultar\\guía rápida};

    % Included use cases (right column)
    \node[ucinc] at ([shift={(3.05, 3.20)}]sys.center) (uceval)
      {Evaluar\\criterios};
    \node[ucinc] at ([shift={(3.05, 1.95)}]sys.center) (ucpers)
      {Persistir\\el caso};

    % Actor associations: bus vertical y líneas horizontales
    \draw[thick] (act.east |- ucmed.west) -- (act.east |- ucguia.west);
    \draw[assoc] (act.east |- ucmed.west) -- (ucmed.west);
    \draw[assoc] (act.east |- ucdx.west) -- (ucdx.west);
    \draw[assoc] (act.east |- uclab.west) -- (uclab.west);
    \draw[assoc] (act.east |- ucver.west) -- (ucver.west);
    \draw[assoc] (act.east |- ucexp.west) -- (ucexp.west);
    \draw[assoc] (act.east |- ucges.west) -- (ucges.west);
    \draw[assoc] (act.east |- ucvis.west) -- (ucvis.west);
    \draw[assoc] (act.east |- ucguia.west) -- (ucguia.west);

    % Include: captura → evaluar / persistir (etiquetas sobre las líneas)
    \coordinate (capjoin) at ([shift={(0.15, 2.58)}]sys.center);
    \draw[inc] (ucmed.east) -- ++(0.45,0) |- (capjoin);
    \draw[inc] (ucdx.east) -- ++(0.45,0) |- (capjoin);
    \draw[inc] (uclab.east) -- ++(0.45,0) |- (capjoin);
    \draw[inc] (capjoin) --
      node[above, font=\scriptsize] {«include»}
      (uceval.west);
    \draw[inc] (capjoin) --
      node[below, font=\scriptsize] {«include»}
      (ucpers.west);
  \end{tikzpicture}%
  }
  \caption{Diagrama de casos de uso. Las líneas continuas asocian al
  actor con cada objetivo. Las discontinuas marcan la relación
  \texttt{include}: evaluar y persistir se ejecutan siempre que cambia
  la captura, sin invocación aparte.}
  \label{fig:casos-uso}
\end{figure}

\section{Análisis}
\label{sec:analisis}

El análisis de dominio distingue el \textbf{caso clínico} que evalúa el
motor del \textbf{catálogo de criterios}, que no forma parte del estado de
sesión. Para interpretar el alcance se separan tres niveles de información:

\begin{itemize}
  \item \textbf{Disponible en la interfaz:} medicaciones, diagnósticos, los
    nueve parámetros de analítica y constantes, y marcadores de pestañas
    revisadas.
  \item \textbf{Disponible solo mediante el JSON del caso:} la ficha
    \texttt{PatientInfo} (incluidos edad y sexo) y los campos heredados de
    dosis y duración de \texttt{Med}; la interfaz no permite capturarlos.
  \item \textbf{Sujeta a revisión manual:} condiciones expresadas solo en
    el texto del aviso, como determinados umbrales de dosis o duración. El
    motor puede mostrar el aviso, pero no comprobar automáticamente esa
    condición con la captura ordinaria.
\end{itemize}

Esta clasificación no demuestra cuántos de los 190 enunciados oficiales
son evaluables en cada nivel. Un \texttt{Crit} se relaciona con el caso por
evaluación: dado un \texttt{PatientCase}, resulta aplicable o no. La
\cref{fig:modelo-dominio} conserva únicamente entidades y cardinalidades;
los atributos y la secuencia se detallan en el~\cref{cap:diseno}.

\begin{figure}[H]
  \centering
  \resizebox{\textwidth}{!}{%
  \begin{tikzpicture}[
    font=\small\sffamily,
    ent/.style={
      draw, thick, rounded corners=2pt, align=center,
      minimum width=2.6cm, minimum height=0.85cm,
      inner sep=3pt, fill=gray!6
    },
    agg/.style={
      draw, thick, rounded corners=2pt, align=center,
      minimum width=2.8cm, minimum height=0.95cm,
      inner sep=3pt, fill=gray!16
    }
  ]
    \node[agg] (pc) {PatientCase};
    \node[ent, above left=1.15cm and 1.15cm of pc] (info) {PatientInfo};
    \node[ent, above right=1.15cm and 1.15cm of pc] (med) {Med};
    \node[ent, below left=1.15cm and 1.15cm of pc] (dx) {Diagnóstico};
    \node[ent, below=1.35cm of pc] (labs) {Labs};
    \node[ent, below right=1.15cm and 1.15cm of pc] (crit) {Crit};

    \draw[thick] (pc) -- (info)
      node[pos=0.12, left, font=\scriptsize] {1}
      node[midway, above, sloped, font=\scriptsize] {ficha}
      node[pos=0.88, right, font=\scriptsize] {0..1};
    \draw[thick] (pc) -- (med)
      node[pos=0.12, right, font=\scriptsize] {1}
      node[midway, above, sloped, font=\scriptsize] {medicaciones}
      node[pos=0.88, left, font=\scriptsize] {0..*};
    \draw[thick] (pc) -- (dx)
      node[pos=0.12, left, font=\scriptsize] {1}
      node[midway, above, sloped, font=\scriptsize] {diagnósticos}
      node[pos=0.88, right, font=\scriptsize] {0..*};
    \draw[thick] (pc) -- (labs)
      node[pos=0.12, left, font=\scriptsize] {1}
      node[midway, left, font=\scriptsize] {analítica}
      node[pos=0.88, left, font=\scriptsize] {0..1};
    \draw[thick, dashed] (crit) -- (pc)
      node[midway, above, sloped, font=\scriptsize] {se evalúa frente a}
      node[pos=0.12, right, font=\scriptsize] {*}
      node[pos=0.88, right, font=\scriptsize] {1};
  \end{tikzpicture}%
  }
  \caption{Entidades del dominio, con el nombre de cada relación sobre la
  línea y las multiplicidades en los extremos. \texttt{Crit} pertenece al
  catálogo; el resto constituye el caso de sesión.}
  \label{fig:modelo-dominio}
\end{figure}
